% TOP_PROBLEM: {"problemId": "problem.combinatorial-boolean-matrix-multiplication-in-near-linear-time", "canonicalTitle": "Combinatorial Boolean Matrix Multiplication in Near-Linear Time", "releaseRank": 480, "primaryDomain": "theoretical_computer_science", "primaryDomainLabel": "Theoretical computer science", "displayStatus": "open", "releaseStatus": "open"}
% !TeX program = lualatex
% Definition and sources only; this file is not an LLM solution attempt.
\documentclass[11pt]{article}
\usepackage[margin=25mm]{geometry}
\usepackage{fontspec}
\setmainfont{Latin Modern Roman}
\usepackage{amsmath,amssymb,mathtools}
\usepackage{unicode-math}
\setmathfont{Latin Modern Math}
\usepackage{xurl,hyperref}
\hypersetup{colorlinks=true,urlcolor=blue}
\setcounter{secnumdepth}{0}
\setlength{\parindent}{0pt}
\setlength{\parskip}{5pt}
\setlength{\emergencystretch}{3em}
\begin{document}
\section*{\texorpdfstring{Combinatorial Boolean Matrix Multiplication in Near-Linear Time}{Combinatorial Boolean Matrix Multiplication in Near-Linear Time}}\label{480}
\subsection{Definitions and mathematical statement}
For Boolean matrices $A,B\in\{0,1\}^{n\times n}$, their Boolean product is
\[
 C_{ij}=\bigvee_{k=1}^n(A_{ik}\wedge B_{kj}).
\]
The target asks for a purely combinatorial algorithm computing all $n^2$ output bits in $O(n^{2+o(1)})$ time. Here the exponent loss tends to zero as $n\to\infty$, so for every fixed ${\varepsilon}>0$ the bound is eventually $O(n^{2+{\varepsilon}})$.

The exclusion of algebraic Strassen-like tensor methods is part of the record. However, purely combinatorial is not itself a machine model or a universally formalized class of algorithms. A fully formal solution must specify allowed operations and justify that it meets the intended exclusion, using the cited computational convention. A generic fast algebraic matrix-multiplication algorithm is not silently counted as a solution to this restricted-method problem.
\par\smallskip\textit{Formulation and concept references:} \href{https://doi.org/10.1109/FOCS.2015.19}{[S1]}.

\subsection{Short English statement}
Can Boolean matrix multiplication be performed in nearly quadratic time using a genuinely combinatorial rather than algebraic tensor-based method?
\subsection{Sources}
\begin{itemize}
\item[S1]H. Yu, FOCS 2015: 109-119, 2015.
\url{https://doi.org/10.1109/FOCS.2015.19}.
\textit{Formulation source.}
\end{itemize}
\end{document}
